\clearpage
\section{프로베니우스 자동동형}
\label{sec:finite-field-frobenius}

\begin{learninggoal}
유한체 확장의 모든 기준체 자동동형을 상대 프로베니우스의 거듭제곱으로 기술한다. 프로베니우스의 위수와 고정체를 계산한다.
\end{learninggoal}

$q=p^r$이라 하고 유한확장
\[
  \F_q\subseteq\F_{q^n}
\]
을 생각하자. 표수 $p$의 프로베니우스 $x\mapsto x^p$를 $r$번 합성하면
\[
  x\longmapsto x^q
\]
를 얻는다.

\begin{definition}[상대 프로베니우스]
사상
\[
  \Frob_q:\F_{q^n}\longrightarrow\F_{q^n},
  \qquad
  a\longmapsto a^q
\]
를 $\F_{q^n}/\F_q$의 \emph{상대 프로베니우스 자동동형}(relative Frobenius automorphism)이라 한다.
\end{definition}

\begin{proposition}
$\Frob_q$는 $\F_q$의 모든 원소를 고정하는 체자동동형이다.
\end{proposition}

\begin{proof}
$q$가 $p$의 거듭제곱이므로
\[
  (a+b)^q=a^q+b^q,
  \qquad
  (ab)^q=a^qb^q.
\]
따라서 $\Frob_q$는 체 준동형이고 단사이다. 정의역이 유한하므로 전사이기도 하다. 또한 $c\in\F_q$이면 $c^q=c$이므로 $\F_q$를 점별로 고정한다.
\end{proof}

\begin{theorem}[유한체 확장의 자동동형군]
$[\F_{q^n}:\F_q]=n$이고
\[
  \Aut_{\F_q}(\F_{q^n})
  =\langle\Frob_q\rangle
\]
이다. 특히 이 자동동형군은 위수 $n$인 순환군이며
\[
  \Aut_{\F_q}(\F_{q^n})
  =\{\id,\Frob_q,\Frob_q^2,\dots,\Frob_q^{n-1}\}.
\]
\end{theorem}

\begin{proof}
모든 $a\in\F_{q^n}$에 대해 $a^{q^n}=a$이므로
\[
  \Frob_q^n=\id.
\]
만약 어떤 $0<d<n$에 대해 $\Frob_q^d=\id$라면 $\F_{q^n}$의 모든 원소가
\[
  X^{q^d}-X
\]
의 근이 된다. 그러나 이 다항식의 차수는 $q^d<q^n$이므로 $q^n$개의 서로 다른 근을 가질 수 없다. 따라서 $\Frob_q$의 위수는 정확히 $n$이다.

한편 $\F_{q^n}/\F_q$는 정규이고 분리인 차수 $n$의 확장이므로 $\F_q$를 고정하는 자동동형은 정확히 $n$개이다. 이미 $\Frob_q$의 거듭제곱으로 $n$개의 서로 다른 자동동형을 얻었으므로 이것들이 전부이다.
\end{proof}

\begin{remark}
다음 장부터 정규이면서 분리인 확장을 갈루아 확장이라 하고 그 자동동형군을 갈루아군이라 부른다. 따라서 위 정리는
\[
  \Gal(\F_{q^n}/\F_q)\cong\Z/n\Z
\]
라는 유한체 갈루아군의 완전한 계산을 미리 제공한다.
\end{remark}

\subsection{프로베니우스의 고정체}

\begin{proposition}
정수 $d\ge0$에 대해 $\Frob_q^d$의 고정체는
\[
  \{a\in\F_{q^n}:a^{q^d}=a\}
  =\F_{q^{\gcd(n,d)}}
\]
이다. 특히 $d=0$일 때에는 고정체를 $\F_{q^n}$으로 해석한다.
\end{proposition}

\begin{proof}
고정원소들은 $X^{q^d}-X$의 근이면서 $\F_{q^n}$에 속하는 원소들이다. 고정된 대수적 폐포 안에서 이는
\[
  \F_{q^d}\cap\F_{q^n}
\]
과 같고, 앞 절의 교집합 공식에 의해 $\F_{q^{\gcd(n,d)}}$이다. $d=0$이면 $\Frob_q^0=\id$이므로 전체 체가 고정된다.
\end{proof}

\begin{corollary}
$\F_{q^n}/\F_q$의 중간체와 $\Aut_{\F_q}(\F_{q^n})$의 부분군은 모두 $n$의 약수로 매개화된다. $d\mid n$에 대해
\[
  \F_{q^d}
\]
는 부분군
\[
  \langle\Frob_q^d\rangle
\]
의 고정체이다.
\end{corollary}

\begin{proof}
$\langle\Frob_q^d\rangle$의 모든 원소를 고정하는 것은 생성원 $\Frob_q^d$를 고정하는 것과 같고, $d\mid n$이면 $\gcd(n,d)=d$이다. 부분체와 부분군의 전체 분류는 각각 약수에 대한 앞의 결과와 순환군의 부분군 분류에서 따른다.
\end{proof}

\subsection{구체적인 프로베니우스 궤도}

\begin{example}[$\F_8/\F_2$]
$\alpha^3+\alpha+1=0$이라 하자. 상대 프로베니우스는
\[
  \Frob_2(a)=a^2
\]
이고
\[
  \alpha\longmapsto\alpha^2
  \longmapsto\alpha^4=\alpha^2+\alpha
  \longmapsto\alpha^8=\alpha
\]
이다. 따라서 자동동형군은 위수 $3$인 순환군이다. $\alpha$의 세 프로베니우스 켤레는 $X^3+X+1$의 세 근이다.
\end{example}

\begin{example}[$\F_9/\F_3$]
$\F_9=\F_3(i)$, $i^2=-1$이라 하자. 상대 프로베니우스는
\[
  a+bi\longmapsto(a+bi)^3=a-bi
\]
이다. 특히 $i\mapsto-i$이고 두 번 적용하면 항등사상이 된다. 따라서
\[
  \Aut_{\F_3}(\F_9)=\{\id,\Frob_3\}\cong\Z/2\Z.
\]
\end{example}

\begin{keyidea}
유한체의 모든 대칭은 하나의 사상에서 나온다.
\[
  a\longmapsto a^q.
\]
확장차수 $n$은 이 상대 프로베니우스를 몇 번 적용해야 항등사상으로 돌아오는지를 나타낸다.
\end{keyidea}

\begin{checkpoint}
\begin{enumerate}[label=(\alph*)]
  \item $\Aut_{\F_5}(\F_{5^6})$의 원소를 모두 기호로 적어라.
  \item $\Frob_2^4$가 $\F_{2^{12}}$에서 고정하는 부분체를 구하라.
  \item $\F_{3^6}/\F_{3^2}$의 상대 프로베니우스와 그 위수를 구하라.
\end{enumerate}
\end{checkpoint}
